% Exact source-aware P21.3 curve-obstruction packet, 25 August 2026.
% No novelty or priority claim is intended.

\paragraph{Source scope.}
The current convention is that on a pure one-dimensional analytic space a
current of bidegree $(0,0)$ acts on compactly supported $(1,1)$-forms.  The
only imported facts are the following standard parts of the controlling
calculi.
\begin{enumerate}
\item Chambert--Loir--Ducros, \emph{Formes diff\'erentielles r\'eelles et
courants sur les espaces de Berkovich}, arXiv:1204.6277v2: proper
tropical pushforward is the transpose of pullback on compactly supported
test forms, preserves order zero and positivity, and commutes with
$\mathrm d'$ and $\mathrm d''$; a balanced weighted polyhedral integration
current is closed.  The converse needed for the particular one-dimensional
fan below is proved directly, rather than imported.  The audited v2 source
archive has SHA--256
\texttt{beca61709a5e35b6f3102948bab004b35c21afdce040ad99d67b15da3de1b11d}.
\item Gubler--K\"unnemann, \emph{A tropical approach to non-Archimedean
Arakelov geometry}, and \emph{Positivity properties of metrics and
delta-forms}: delta-forms are compatible under common refinements of
tropical charts; restriction of a positive closed delta-current to smooth
test forms is a positive closed CLD current; and the tropical
Poincar\'e--Lelong/corner formula computes the delta-current of a real PL
function.
\end{enumerate}
No pullback of an arbitrary current is used.  Every current below is pushed
forward by transposition against target test forms.

\begin{proof}
Pairwise distinct reductions give
\[
 |a-b|=|a-c|=\cdots=|c-d|=1.
\tag{1}
\]
The convex hull in the Berkovich projective line of the four punctures is
therefore the star with central vertex the unit Gauss point $q$ and four
unbounded open edges, one toward each puncture.  Both $f$ and $g$ are
analytic units on $X$.  On the four edges their logarithmic slopes are
\[
\begin{array}{c|rrrr}
 &a&b&c&d\\ \hline
s=-\log|f|& 1&-1&0&0\\
t=-\log|g|& 0&0&1&-1.
\end{array}
\tag{2}
\]
Outside this star both norms are constant on the fibres of the retraction.
Thus $\tau(X)$ is the full cross and its four primitive ray vectors are the
ones in the statement.  The analytic multiplicity and every tangent-lattice
index are one, so the fundamental weights are one.

For completeness, properness is algebraic rather than an assumption hidden
in the picture.  The M\"obius coordinate $x=f$ identifies
$\mathbf P^1\setminus\{a,b\}$ with $\mathbf G_m$.  If
$\alpha=f(c)$ and $\beta=f(d)$, then $\alpha,\beta\in K^\times$ are
distinct and there is $C\in K^\times$ such that
\[
                         g=C\frac{x-\alpha}{x-\beta}.
\tag{3}
\]
Hence $X$ is the closed curve
$g(x-\beta)=C(x-\alpha)$ in $\mathbf G_m^2$: at $x=\alpha$ or
$x=\beta$ the equation has no solution in the torus.  Tropicalization of a
closed subvariety of a torus is proper, which proves the claim for $\tau$.

The ambient functions $2|s|$ and $|s|+|t|$ are continuous convex
functions.  The smooth convex functions
\[
 2\sqrt{s^2+\varepsilon},\qquad
 \sqrt{s^2+\varepsilon}+\sqrt{t^2+\varepsilon}
\]
decrease locally uniformly to them as $\varepsilon\downarrow0$.  Pullback
by the moment map therefore proves that $U,V$ are continuous psh, including
the CLD approximation requirement.

The one-dimensional tropical Poincar\'e--Lelong formula says that the atom
of the curvature of a PL function at an interior vertex is the sum of its
outgoing slopes, counted with the canonical weights.  By (2), the outgoing
slopes of $U$ at $q$ are $2,2,0,0$, while those of $V$ are $1,1,1,1$.
Both sums are four.  The functions are affine on every open star edge and
constant on every complementary retraction fibre; the punctures are not
points of $X$.  Hence there is no further curvature term and
\[
             \mathrm d'\mathrm d''[U]
             =\mathrm d'\mathrm d''[V]=4\delta_q.
\tag{4}
\]
The same corner-cycle computation is an equality of delta-currents.

It remains to test the proposed current against all refinements, rather
than only against the coarse chart $(f,g)$.  Adjoin the analytic unit
\[
                         k=\frac{z-a}{z-c},
                 \qquad r=-\log|k|,
\]
and put $\widehat\tau=(s,t,r)$.  In the coordinate $x=f$ there is
$D\in K^\times$ such that
\[
                         k=D\frac{x}{x-\alpha}.
\tag{5}
\]
Equations (3) and (5) exhibit $X$ as the closed curve
\[
 g(x-\beta)=C(x-\alpha),\qquad k(x-\alpha)=Dx
\]
in $\mathbf G_m^3$.  Thus $\widehat\tau$ is again proper.  Its four
primitive ray vectors, in the order $a,b,c,d$, are
\[
 v_a=(1,0,1),\quad v_b=(-1,0,0),\quad
 v_c=(0,1,-1),\quad v_d=(0,-1,0).
\tag{6}
\]
Every vector in (6) is primitive, and projection to the first two
coordinates maps its ray isomorphically, with lattice index one, onto the
corresponding ray of the coarse cross.

Suppose that $T$ is a positive closed order-zero CLD current of bidegree
$(0,0)$ and write $S=\widehat\tau_*T$.  Proper pushforward makes $S$ a
positive closed order-zero supercurrent supported on the fan (6).  On the
relative interior of one ray choose integral affine coordinates
$y_1,y_2,y_3$, with $y_1$ increasing in the primitive ray direction and
the ray given by $y_2=y_3=0$.  An order-zero positive supercurrent of
bidimension $(1,1)$ is represented there by a positive-semidefinite
matrix of Radon measures.  The equations $\mathrm d'S=\mathrm d''S=0$,
tested against forms whose coefficients have arbitrary normal first jets,
kill every row and column except the $y_1y_1$ entry.  The remaining
equation is the distributional identity $\partial_{y_1}\mu=0$.  Hence
$\mu=w_e\,\mathrm dy_1$ for a constant $w_e\geq0$ on that ray.

An order-zero bidimension-$(1,1)$ remainder at the vertex would have the
form $A\delta_0$ for a positive-semidefinite matrix $A$.  Its differential
contains first derivatives of $\delta_0$, whereas the boundaries of the
four ray currents contain only multiples of $\delta_0$.  Testing against
forms with prescribed first jet at the origin therefore gives $A=0$.
Finally, testing closedness against forms constant near the vertex gives
the boundary equation
\[
                         \sum_{e=a,b,c,d}w_ev_e=0.
\tag{7}
\]
This is the elementary order-zero support calculation behind the necessity
of balancing in this one-dimensional case.  The weights in (7) agree with
the weights of $\tau_*T$ because
the projection has one ray over each ray and all four lattice indices are
one.

Substituting (6) in (7) gives
\[
 (w_a-w_b,\ w_c-w_d,\ w_a-w_c)=(0,0,0),
\]
and therefore
\[
                         w_a=w_b=w_c=w_d.
\tag{8}
\]
In particular $(w_a,w_b,w_c,w_d)=(2,2,1,1)$ is impossible; its refined
balancing defect is exactly $(0,0,1)$.

Finally let $T_\delta$ be a positive closed delta-current with the asserted
visible ray weights.  Restriction to smooth test forms commutes with both
differentials and preserves positivity.  Its CLD restriction is therefore
a positive closed current to which (8) applies.  Since one-dimensional ray
weights are detected by smooth compactly supported top-degree forms on the
ray interiors, restriction does not discard any of the four weights.  The
same contradiction follows.  This proves every assertion.
\end{proof}

\paragraph{Scope firewall.}
The result rejects the concrete four-punctured smooth-curve realization.
It does not claim that every positive closed bidegree-$(0,0)$ current on an
arbitrary connected analytic curve is globally a scalar multiple of the
fundamental current, and it does not decide the residual problem for
higher-dimensional analytic CLD currents or for non-PL potentials.
